\documentclass[10pt,a4paper]{amsart}
\usepackage{amsmath,amssymb,amsthm,graphicx}
% Retained minimal preamble for the Stacks Project chapter more-algebra.tex (More on Algebra).
% The chapter file pinned-original-source.tex carries only \input{preamble} on its first line, so no
% preamble text lives inside the chapter itself.  The Stacks Project uses one single shared preamble for
% every chapter; the snapshot stages it identically as sources/stacks-chow/preamble.tex and
% sources/stacks-etale-cohomology/preamble.tex.  The block below is copied mechanically, line for line,
% from sources/stacks-chow/preamble.tex lines 14-18 (the Xy-pic setup) and 154-217 (the native theorem
% declarations together with the mathematical macros).  Nothing was added, removed or redefined; the
% wrapper class and the amsmath/amssymb/amsthm/graphicx base packages are supplied by the assembler.
% Chapter: more-algebra.tex -- More on Algebra.
% For commutative diagrams we use Xy-pic
\usepackage[all]{xy}

% We use 2cell for 2-commutative diagrams.
\xyoption{2cell}

% Theorem environments.
%
\theoremstyle{plain}
\newtheorem{theorem}[subsection]{Theorem}
\newtheorem{proposition}[subsection]{Proposition}
\newtheorem{lemma}[subsection]{Lemma}

\theoremstyle{definition}
\newtheorem{definition}[subsection]{Definition}
\newtheorem{example}[subsection]{Example}
\newtheorem{exercise}[subsection]{Exercise}
\newtheorem{situation}[subsection]{Situation}

\theoremstyle{remark}
\newtheorem{remark}[subsection]{Remark}
\newtheorem{remarks}[subsection]{Remarks}

\numberwithin{equation}{subsection}

% Macros
%
\def\lim{\mathop{\mathrm{lim}}\nolimits}
\def\colim{\mathop{\mathrm{colim}}\nolimits}
\def\Spec{\mathop{\mathrm{Spec}}}
\def\Hom{\mathop{\mathrm{Hom}}\nolimits}
\def\Ext{\mathop{\mathrm{Ext}}\nolimits}
\def\SheafHom{\mathop{\mathcal{H}\!\mathit{om}}\nolimits}
\def\SheafExt{\mathop{\mathcal{E}\!\mathit{xt}}\nolimits}
\def\Sch{\mathit{Sch}}
\def\Mor{\mathop{\mathrm{Mor}}\nolimits}
\def\Ob{\mathop{\mathrm{Ob}}\nolimits}
\def\Sh{\mathop{\mathit{Sh}}\nolimits}
\def\NL{\mathop{N\!L}\nolimits}
\def\CH{\mathop{\mathrm{CH}}\nolimits}
\def\proetale{{pro\text{-}\acute{e}tale}}
\def\etale{{\acute{e}tale}}
\def\QCoh{\mathit{QCoh}}
\def\Ker{\mathop{\mathrm{Ker}}}
\def\Im{\mathop{\mathrm{Im}}}
\def\Coker{\mathop{\mathrm{Coker}}}
\def\Coim{\mathop{\mathrm{Coim}}}

% Boxtimes
%
\DeclareMathSymbol{\boxtimes}{\mathbin}{AMSa}{"02}

%
% Macros for moduli stacks/spaces
%
\def\QCohstack{\mathcal{QC}\!\mathit{oh}}
\def\Cohstack{\mathcal{C}\!\mathit{oh}}
\def\Spacesstack{\mathcal{S}\!\mathit{paces}}
\def\Quotfunctor{\mathrm{Quot}}
\def\Hilbfunctor{\mathrm{Hilb}}
\def\Curvesstack{\mathcal{C}\!\mathit{urves}}
\def\Polarizedstack{\mathcal{P}\!\mathit{olarized}}
\def\Complexesstack{\mathcal{C}\!\mathit{omplexes}}
% \Pic is the operator that assigns to X its picard group, usage \Pic(X)
% \Picardstack_{X/B} denotes the Picard stack of X over B
% \Picardfunctor_{X/B} denotes the Picard functor of X over B
\def\Pic{\mathop{\mathrm{Pic}}\nolimits}
\def\Picardstack{\mathcal{P}\!\mathit{ic}}
\def\Picardfunctor{\mathrm{Pic}}
\def\Deformationcategory{\mathcal{D}\!\mathit{ef}}
\usepackage{cleveref}
\crefname{originalremark}{original remark}{original remarks}
\Crefname{originalremark}{Original remark}{Original remarks}
\makeatletter
\newlabel{example-eta-not-distinguished}{{0GSN}{source}{}{sourcepointer-1}{}}
\newlabel{example-eta-not-distinguished@cref}{{[originalremark][0][]0GSN}{[1][0][]source}{}{}{}}
\makeatother
\crefname{originallemma}{original lemma}{original lemmas}
\Crefname{originallemma}{Original lemma}{Original lemmas}
\makeatletter
\newlabel{derived-lemma-localization-subcategory}{{0GSL}{source}{}{sourcepointer-2}{}}
\newlabel{derived-lemma-localization-subcategory@cref}{{[originallemma][0][]0GSL}{[1][0][]source}{}{}{}}
\makeatother
\crefname{originallemma}{original lemma}{original lemmas}
\Crefname{originallemma}{Original lemma}{Original lemmas}
\makeatletter
\newlabel{lemma-K-flat-resolution}{{06Y4}{source}{}{sourcepointer-3}{}}
\newlabel{lemma-K-flat-resolution@cref}{{[originallemma][0][]06Y4}{[1][0][]source}{}{}{}}
\makeatother
\crefname{originallemma}{original lemma}{original lemmas}
\Crefname{originallemma}{Original lemma}{Original lemmas}
\makeatletter
\newlabel{lemma-eta-qis}{{0F7Q}{source}{}{sourcepointer-4}{}}
\newlabel{lemma-eta-qis@cref}{{[originallemma][0][]0F7Q}{[1][0][]source}{}{}{}}
\makeatother
\crefname{originallemma}{original lemma}{original lemmas}
\Crefname{originallemma}{Original lemma}{Original lemmas}
\makeatletter
\newlabel{categories-lemma-properties-left-localization}{{04VG}{source}{}{sourcepointer-5}{}}
\newlabel{categories-lemma-properties-left-localization@cref}{{[originallemma][0][]04VG}{[1][0][]source}{}{}{}}
\makeatother
\title{The shifted-complex construction descends to an additive functor on the derived category for a nonzerodivisor}
\author{The Stacks Project Authors}
\date{stacks-project:more-algebra.tex:local-snapshot-20261012; GNU FDL 1.2 or later}
\begin{document}
\maketitle
\setcounter{subsection}{1}
\noindent\textit{Editorial note.} All delivered mathematical text is the author's own, copied from the pinned Stacks Project chapter file more-algebra.tex; no statement or proof step was rewritten, reordered or removed. The counted statement is the labeled lemma environment at source lines 27531-27540 (canonical tag 0F7R) and the counted proof is the proof environment at source lines 27542-27576, which immediately follows it and is its own proof.\par\medskip
\noindent\textit{Editorial note.} The statement carries one footnote warning that the functor is not exact and pointing to an example; the footnote is part of the author's own statement inside the counted span and its reference is carried as a navigation stub.\par\medskip
\noindent\textit{Editorial note.} The proof refers to a localization criterion and a formal property of localizations from other chapters and to two earlier results of this chapter; all five are navigation stubs with their roles recorded in the reasons.\par\medskip
\noindent\textit{Editorial note.} The chapter file pinned-original-source.tex carries only an input of the shared preamble on its first line, so no preamble text lives inside the chapter. The retained file sources/stacks-more-algebra/retained-preamble.tex is a mechanical, line-for-line copy of the project's single shared preamble as staged in the snapshot (Xy-pic setup plus the native theorem declarations and mathematical macros). Nothing was added, removed or redefined, and the excerpt serves only as preamble, never as a span coordinate.\par\medskip
\noindent\textit{Editorial note.} The source URL is the source-level Stacks Project URL, as required when more than one card is produced from one source; the per-claim tag is carried by claim\_key and by this note. The counted statement and proof cite no bibliography entry, so this card needs no bibliography.\par\medskip
\noindent\textit{Allowed standard prerequisites (editorial).}\begin{itemize}
\item The homotopy category of a full subcategory closed under the relevant operations embeds as a full triangulated subcategory, and a localization criterion characterizes when the localized category is equivalent to the target. invoked at proof lines 27543-27556. Reduces the construction of the functor on the derived category to the single requirement that every complex admits a quasi-isomorphism from a complex in the subcategory.
\item Every complex of modules admits a quasi-isomorphism from a K-flat complex of flat modules. invoked at proof lines 27560-27563. Supplies the required resolution, since flat modules are torsion free with respect to a nonzerodivisor.
\item A functor on the homotopy category of the subcategory that sends quasi-isomorphisms to quasi-isomorphisms factors through the localization. invoked at proof lines 27566-27575. Produces the additive functor on the derived category whose value on a complex of torsion free modules is the shifted complex.
\end{itemize}\medskip
\ifdefined\hypertarget\hypertarget{sourcepointer-1}{}\fi
\noindent\textit{Source reference.} Original remark 0GSN: sources/stacks-more-algebra/pinned-original-source.tex, source lines 27445--27445. The reference is cited in a footnote of the statement to warn that the functor does not preserve distinguished triangles; it is a pointer to an example and carries no part of the construction that is proved inside the counted span.\par\medskip
\ifdefined\hypertarget\hypertarget{sourcepointer-2}{}\fi
\noindent\textit{Source reference.} Original lemma 0GSL: sources/stacks-chow/tags.txt, source lines 20936--20936. The reference is the localization criterion applied at proof lines 27550-27556 to reduce the construction to the existence of a torsion free resolution; the resolution argument that discharges it is inside the counted span.\par\medskip
\ifdefined\hypertarget\hypertarget{sourcepointer-3}{}\fi
\noindent\textit{Source reference.} Original lemma 06Y4: sources/stacks-more-algebra/pinned-original-source.tex, source lines 14380--14380. The reference supplies the quasi-isomorphism from a complex of flat modules used at proof lines 27560-27563; flat modules are torsion free with respect to the nonzerodivisor, which is the observation the proof makes.\par\medskip
\ifdefined\hypertarget\hypertarget{sourcepointer-4}{}\fi
\noindent\textit{Source reference.} Original lemma 0F7Q: sources/stacks-more-algebra/pinned-original-source.tex, source lines 27518--27518. The reference is applied at proof lines 27572-27573 to see that the shifted construction sends quasi-isomorphisms to quasi-isomorphisms, which is the hypothesis for factoring through the localization; the factorization argument is inside the counted span.\par\medskip
\ifdefined\hypertarget\hypertarget{sourcepointer-5}{}\fi
\noindent\textit{Source reference.} Original lemma 04VG: sources/stacks-chow/tags.txt, source lines 6324--6324. The reference is the formal property of localizations that lets a functor sending the inverted arrows to isomorphisms factor through the localized category, invoked at proof lines 27573-27575; it is a pointer and the verification of its hypothesis is inside the counted span.\par\medskip
\begin{lemma}
\label{lemma-Leta}
Let $A$ be a ring and let $f \in A$ be a nonzerodivisor. There is an additive
functor\footnote{Beware that this functor isn't exact, i.e.,
does not transform distinguished triangles into distinguished triangles.
See Example \ref{example-eta-not-distinguished}.}
$L\eta_f : D(A) \to D(A)$ such that if $M \in D(A)$ is
represented by a complex $M^\bullet$ of $f$-torsion free $A$-modules,
then $L\eta_fM = \eta_fM^\bullet$ and similarly for morphisms.
\end{lemma}

\begin{proof}
Denote $\mathcal{T} \subset \text{Mod}_A$ the full subcategory
of $f$-torsion free $A$-modules. We have a corresponding inclusion
$$
K(\mathcal{T}) \quad\subset\quad K(\text{Mod}_A) = K(A)
$$
of $K(\mathcal{T})$ as a full triangulated subcategory of $K(A)$.
Let $S \subset \text{Arrows}(K(\mathcal{T}))$ be the quasi-isomorphisms.
We will apply
Derived Categories, Lemma \ref{derived-lemma-localization-subcategory}
to show that the map
$$
S^{-1}K(\mathcal{T}) \longrightarrow D(A)
$$
is an equivalence of triangulated categories. The lemma shows that
it suffices to prove: given a complex $M^\bullet$ of $A$-modules,
there exists a quasi-isomorphism $K^\bullet \to M^\bullet$ with $K^\bullet$
a complex of $f$-torsion free modules.
By Lemma \ref{lemma-K-flat-resolution} we can find a quasi-isomorphism
$K^\bullet \to M^\bullet$ such that the complex $K^\bullet$
is K-flat (we won't use this) and consists of flat $A$-modules $K^i$.
In particular, $f$ is a nonzerodivisor on $K^i$ for all $i$ as desired.

\medskip\noindent
With these preliminaries out of the way we can define $L\eta_f$.
Namely, by the discussion at the start of this section
we have already a well defined functor
$$
K(\mathcal{T}) \xrightarrow{\eta_f} K(\mathcal{T}) \to K(A) \to D(A)
$$
which according to Lemma \ref{lemma-eta-qis} sends quasi-isomorphisms
to quasi-isomorphisms. Hence this functor factors over
$S^{-1}K(\mathcal{T}) = D(A)$ by 
Categories, Lemma \ref{categories-lemma-properties-left-localization}.
\end{proof}

\end{document}
